\documentclass[12pt,letterpaper]{article}

% Georgia College Economics, "Style Guide for Papers" (user-supplied).
% Page setup: US Letter, one-inch margins, double-spaced ragged-right
% Times-family 12pt body text; tables/figures use 11pt.
\usepackage[margin=1in]{geometry}
\usepackage{iftex}
\ifPDFTeX
  % pdfLaTeX cannot embed system Times New Roman directly; NewTX is
  % metrically compatible. Use XeLaTeX for the exact system font.
  \usepackage{newtxtext}
  \usepackage{newtxmath}
\else
  \usepackage{fontspec}
  \IfFontExistsTF{Times New Roman}
    {\setmainfont{Times New Roman}}
    {\IfFontExistsTF{Tinos}{\setmainfont{Tinos}}{\setmainfont{Latin Modern Roman}}}
  \usepackage{unicode-math}
  \IfFontExistsTF{STIX Math}
    {\setmathfont{STIX Math}}
    {\setmathfont{Latin Modern Math}}
\fi
\usepackage{amsmath}
\usepackage{setspace}
\usepackage{ragged2e}
\usepackage{indentfirst}
\usepackage{fancyhdr}
\usepackage{titlesec}
\usepackage{caption}
\usepackage[authoryear,round]{natbib}
\usepackage{etoolbox}

% Paragraph indents persist even though the right margin is ragged.
\RaggedRight
\setlength{\parindent}{0.5in}
\setlength{\parskip}{0pt}
\setlength{\emergencystretch}{3em}
\doublespacing

% Make section headings conform to the paper's 12-point typography.
\titleformat{\section}{\normalfont\bfseries\fontsize{12}{14}\selectfont}
  {\thesection}{1em}{}
\titleformat{\subsection}{\normalfont\bfseries\fontsize{12}{14}\selectfont}
  {\thesubsection}{1em}{}
\titleformat{\subsubsection}{\normalfont\bfseries\fontsize{12}{14}\selectfont}
  {\thesubsubsection}{1em}{}

% Bottom-right page numbers; title page deliberately has none.
\pagestyle{fancy}
\fancyhf{}
\renewcommand{\headrulewidth}{0pt}
\fancyfoot[R]{\normalfont\normalsize\thepage}
\setlength{\headheight}{14pt}

% Keep floats readable: 11pt (the \small size for a 12pt document).
% Narrow column padding preserves text size in dense econometric tables.
\AtBeginEnvironment{table}{%
  \small\singlespacing
  \setlength{\tabcolsep}{2pt}%
  \renewcommand{\arraystretch}{1.2}%
}
\AtBeginEnvironment{figure}{\small\singlespacing}
\captionsetup{font=small,labelfont=bf,labelsep=period,
  justification=raggedright,singlelinecheck=false,skip=11pt}
\captionsetup[table]{position=top}
\captionsetup[figure]{position=bottom}

% Equations receive centered displays and bold numbers at the right margin.
\makeatletter
\renewcommand{\tagform@}[1]{%
  \maketag@@@{\bfseries(\ignorespaces#1\unskip\@@italiccorr)}%
}
\makeatother

% Author-year citations and an alphabetical, unnumbered, single-spaced
% bibliography with hanging indents, on a new numbered page.
\bibpunct{(}{)}{;}{a}{,}{,}
\renewcommand{\refname}{Bibliography}
\renewcommand{\bibfont}{\normalsize}
\setlength{\bibhang}{0.5in}

% Title-page metadata: fill these in if this research paper is being
% submitted for a course. Empty fields are intentionally omitted.
\newcommand{\PaperTitle}{Regime-Dependent Predictability of Stock Returns: Evidence from Machine Learning and Econometric Models}
\newcommand{\PaperAuthor}{Surya Pratap Singh Suryavanshi}
\newcommand{\PaperCourse}{}
\newcommand{\PaperSchedule}{}
\newcommand{\PaperProfessor}{}
\newcommand{\PaperDate}{\today}
\newcommand{\PaperDraftLabel}{} % E.g., First Draft

\begin{document}
\begin{titlepage}
  \thispagestyle{empty}
  \begin{singlespace}
    \centering
    \vspace*{1.55in}
    {\bfseries\PaperTitle\par}
    \ifx\PaperDraftLabel\empty\else\PaperDraftLabel\par\fi
    \vspace{2\baselineskip}
    \PaperAuthor\par
    \ifx\PaperCourse\empty\else\PaperCourse\par\fi
    \ifx\PaperSchedule\empty\else\PaperSchedule\par\fi
    \ifx\PaperProfessor\empty\else\PaperProfessor\par\fi
    \PaperDate\par
  \end{singlespace}
\end{titlepage}

% The title page counts as page 1, but shows no page number.
\setcounter{page}{2}
\begin{center}
  \bfseries\PaperTitle
\end{center}
\vspace{\baselineskip}

\begin{abstract}
This research studies whether the predictability of next-month U.S. common-stock returns varies across high- and low-volatility market regimes and whether nonlinear machine-learning methods capture this variation differently from a regularized linear benchmark. The empirical pipeline uses monthly CRSP stock data, a VIX-based regime classification constructed from an expanding historical median, and strictly time-respecting out-of-sample estimation with Elastic Net and XGBoost. The current preliminary specification uses a computationally reduced fixed large-cap universe and a partial predictor set consisting of size and momentum; book-to-market is added in the subsequent Compustat/CCM data-construction stage. Over 83 out-of-sample formation months from January 2019 through November 2025, the tuned Elastic Net collapses to an intercept-only cross-sectional forecast and therefore provides no usable stock-ranking signal. XGBoost retains cross-sectional forecast dispersion, but its mean monthly Spearman rank correlation with realized next-month returns is slightly negative. Equal-weighted XGBoost D10-minus-D1 returns are close to zero overall, whereas the value-weighted spread averages approximately 0.71\% per month and is concentrated in low-VIX months, where it averages approximately 1.43\% per month; the corresponding high-VIX value-weighted spread is essentially zero. These findings are descriptive and preliminary rather than formal evidence of statistical or net-of-cost profitability. They motivate the next stages of the research, including the complete Compustat/CCM predictor set, formal inference, turnover measurement, and transaction-cost robustness.
\end{abstract}

\section{Introduction}

The predictability of stock returns need not be constant across market environments. Periods of elevated market volatility can coincide with changes in risk premia, investor behavior, information processing, liquidity, and the cross-sectional dispersion of returns. This research therefore asks whether next-month stock-return predictability differs between high- and low-volatility regimes and whether a nonlinear machine-learning model captures regime-dependent predictive structure differently from a regularized linear model.

The production empirical design compares Elastic Net and XGBoost using the same three-predictor information set and the same strictly out-of-sample evaluation structure. At formation month $t$, log market equity, book-to-market, momentum 12--2, and the volatility regime are constructed using information available through $t$, while the prediction target is the realized stock return in calendar month $t+1$. Models are estimated only on historical observations whose outcomes are observable by the formation date.

The production panel combines CRSP monthly stock returns with Compustat annual accounting information through the CRSP/Compustat Merged database. Annual book-to-market is assigned using a conservative June availability convention and prior-December market equity. Volatility regimes are based on VIX relative to its expanding historical median, preventing future volatility information from entering the state classification.

Hyperparameters are selected once using only the 1990--1999 pre-OOS period with forward-chaining validation and are frozen thereafter. The OOS evaluation begins in January 2000 and extends through November 2025. This design separates historical model selection from subsequent forecast evaluation and avoids random train/test splits that would violate the temporal structure of the problem.

The production results reveal an important distinction between unconditional forecast accuracy and cross-sectional economic value. Both models have OOS $R^2$ values close to zero relative to a pooled expanding historical-mean benchmark. However, XGBoost displays statistically positive cross-sectional rank information in LOW-VIX months, and equal-weighted prediction-sorted long--short portfolios retain positive returns after the frozen transaction-cost adjustment. The same signal is weak in HIGH-VIX months and disappears under value weighting. At the same time, direct LOW-minus-HIGH tests do not reach conventional significance. The evidence therefore points to predictive content that is concentrated in low-volatility states without establishing a sharp statistically significant break between regimes.

\section{Related Literature, Contribution, and Research Hypotheses}
\label{sec:related-literature}

This study lies at the intersection of cross-sectional empirical asset pricing,
nonlinear statistical forecasting, time-varying market conditions, and the
economics of implementing return predictions. These literatures overlap, but
they address different questions. Evidence that a model predicts the
\emph{aggregate} equity premium is not, by itself, evidence that it ranks
\emph{individual stocks} accurately. Similarly, statistical accuracy,
cross-sectional rank correlation, gross long--short returns, and investable
profitability are separate objects. The contribution of the present study must
be evaluated against all four distinctions.

\subsection{Machine learning and cross-sectional asset pricing}
\label{subsec:literature-ml}

The empirical asset-pricing literature has long studied how the observable
characteristics of securities relate to their subsequent returns.
\cite{FamaFrench1992} identify important cross-sectional associations involving
size and book-to-market, while \cite{JegadeeshTitman1993} document the
profitability of intermediate-horizon momentum strategies.
The common-risk-factor formulation of \cite{FamaFrench1993} motivates a
distinction between a characteristic's predictive association and its relation
to systematic risk. Subsequent research on gross profitability and investment
illustrates the limitations of a specification containing size, value and
momentum alone \cite{NovyMarx2013,FamaFrench2015}. The three characteristics
used here are therefore interpretable and historically established, but should
not be mistaken for a complete information set.

The introduction of flexible prediction methods has shifted attention from
individual characteristic slopes to high-dimensional and potentially nonlinear
maps from firm information to expected returns. \cite{GuKellyXiu2020} provide
a broad comparison of machine-learning methods for empirical asset pricing and
show the importance of allowing nonlinear interactions among predictors.
\cite{FreybergerNeuhierlWeber2020} study characteristic selection and
nonlinearity in a nonparametric framework, while \cite{KozakNagelSantosh2020}
illustrate the value of economically motivated shrinkage in the cross-section.
\cite{KellyPruittSu2019} model characteristics as conditioning information
for time-varying risk exposures. These studies demonstrate why a
regularized linear predictor and a nonlinear boosted-tree model provide
informative benchmarks, but they also establish that neither regularization
nor nonlinear characteristic modeling is a new methodological invention.
In particular, Gu, Kelly and Xiu explicitly include a three-variable
ordinary-least-squares comparator based on size, book-to-market and
momentum, and compare it with regularized linear and gradient-boosted
models. Their approximately 1957--2016 monthly stock panel uses an
1987--2016 out-of-sample period with annual model refits, as opposed
to the monthly expanding refits of the present project. Thus, using
these three characteristics alone cannot be represented as an original
selection of explanatory variables.

Further evidence on flexible factor representations is provided by
\cite{ChenPelgerZhu2024}. The critical survey of
\cite{Bagnara2024} highlights the range of regularization, dimension
reduction, tree and neural-network approaches already applied to asset
pricing. Thus, comparing Elastic Net with XGBoost should be framed as a
controlled empirical comparison of two established forecasting families,
not as the introduction of a new class of forecasting model.

The present design deliberately keeps the characteristic set small.
The frozen predictors are log market equity, book-to-market, and
momentum 12--2, and the same sample and formation-time information set are
used for both models. This restriction makes the model contrast interpretable:
differences between the two forecasting methods cannot be attributed simply
to one model receiving additional firm characteristics. However, it also
means that any apparent nonlinear advantage must be interpreted relative
to a sparse baseline; important known predictors and stronger simple
characteristic benchmarks are reserved for explicit robustness analyses.

\subsection{Volatility, market states, and conditional predictability}
\label{subsec:literature-volatility}

Market volatility summarizes more than a single structural mechanism.
\cite{BekaertHoerova2014} show that the VIX contains information connected
to expected variance and variance risk premia. The relation between
variance risk premia and subsequent aggregate stock-market returns is
studied by \cite{BollerslevTauchenZhou2009}. These contributions help
motivate a VIX-based conditioning variable, but their aggregate-return
outcomes are not identical to individual-stock return rankings.
A VIX state is also observational: it co-moves with risk perceptions,
funding conditions, liquidity and the business cycle, and it is not an
exogenously assigned intervention.

Studies of volatility-managed factor portfolios illustrate that an investor's
exposure to time-varying volatility may matter even if unconditional factor
relationships remain the same \cite{MoreiraMuir2017}.
Momentum provides a particularly relevant example because it belongs to
the current study's frozen characteristic set:
\cite{DanielMoskowitz2016} examine momentum crashes associated with
panic-like market conditions, and \cite{BarrosoSantaClara2015}
investigate changing momentum risk. A high-volatility forecast-performance
difference could therefore reflect an interaction between VIX states and
momentum exposure rather than a general property of nonlinear machine
learning.

The study uses two historically defined, but conceptually distinct,
classifications. The original binary LOW/HIGH scheme compares formation-month
VIX with its expanding historical median. The additional LOW/MIDDLE/HIGH
evaluation uses expanding historical VIX terciles. Both are calculated using
information available at portfolio formation, and neither involves fitting
separate models in each regime. Their thresholds are reproducible empirical
conventions rather than established structural break points. Consequently,
alternative VIX thresholds and realized-volatility measures are necessary
robustness checks before drawing claims about the market-state nature of
predictability.

\subsection{Conditional nonlinear predictability and closest competing studies}
\label{subsec:literature-closest}

A major originality challenge is that volatility-dependent machine-learning
returns and implementation constraints have already been studied directly.
\cite{AvramovChengMetzker2023} find that investment returns from deep-learning
signals are concentrated in difficult-to-arbitrage stocks and deteriorate
when microcapitalization stocks, distressed firms, or high-volatility
episodes are excluded; turnover and trading costs further reduce
profitability. This substantially overlaps the motivation of the present
analysis. It would therefore be incorrect to claim that linking machine
learning, market volatility, microcaps and transaction costs is itself
unprecedented.

The current empirical design differs in several more specific dimensions.
It evaluates historically frozen forecasts from \emph{two} model classes on
the same three-characteristic information set; partitions the same monthly
forecast record using historically observable binary and three-state VIX
classifications; distinguishes stock-level forecast loss from monthly
cross-sectional rank information; tests \emph{direct} differences between
regime-conditional means with heteroskedasticity- and
autocorrelation-consistent inference; and explicitly compares gross and
cost-adjusted equal- and value-weighted decile-spread portfolios.
A recent non-refereed study by \cite{FangSlepaczuk2026} likewise
combines regime-aware volatility forecasting, XGBoost return prediction,
walk-forward validation and transaction-cost discussion. Its Chinese
high-frequency index setting differs from the present monthly
individual-stock cross-section, but the overlapping research question
further limits broad claims of priority.

Those design features can improve interpretability. Nevertheless, assembling
existing methods into one pipeline is not by itself proof of a sufficiently
distinct journal contribution. That determination requires a full-text
comparison with competing studies and further empirical evidence that resolves
a substantive question not already answered.

The literature also warns against attributing all state dependence to
the model architecture. \cite{KellyPruittSu2019} provide a risk-exposure
interpretation of time-varying characteristic relations, while
\cite{DanielMoskowitz2016} highlight the potentially nonlinear downside
risk embedded in momentum. A model may rank stocks differently because
its characteristic responses are nonlinear, because the composition of the
eligible cross-section changes, or because macroeconomic exposures vary by
state. These explanations should be separated as far as observable data
permit, but the current design cannot establish a causal channel.

\subsection{Statistical evaluation, economic value, and investment frictions}
\label{subsec:literature-evaluation}

The evaluation of forecasts requires matching the statistic to the
research question. Mean squared forecasting error, out-of-sample
$R^2$, a monthly Spearman rank correlation, and a realized decile spread
capture different properties of a prediction. Evidence of modest
out-of-sample improvement for aggregate equity-premium forecasts
\cite{CampbellThompson2008,WelchGoyal2008,RapachStraussZhou2010}
provides useful caution about benchmarks and stability, but should not
be presented as direct evidence about the cross-section of individual
securities.

Comparing two forecast models on a common set of realized outcomes
also creates dependence. \cite{DieboldMariano1995} and
\cite{GiacominiWhite2006} provide influential approaches to comparing
predictive accuracy and conditional predictive ability, respectively,
under their stated sampling and estimation assumptions. Those results do
not license treating hundreds of thousands of stock--month residuals
as independent test observations. In this paper the formation month is
the common decision and conditioning unit. Additional model-comparison
tests will therefore form aligned monthly loss or rank-difference series
before applying an appropriate dependent-observation inference procedure.

For the existing monthly mean-return, rank-IC and regime-difference analyses,
the reported framework uses the heteroskedasticity- and
autocorrelation-consistent covariance estimator associated with
\cite{NeweyWest1987}. Direct LOW/MIDDLE/HIGH pairwise contrasts
are evaluated with the family-wise step-down adjustment of
\cite{Holm1979}; the false-discovery-rate procedure of
\cite{BenjaminiHochberg1995} is a separately identified secondary
diagnostic. This distinction is important: a positive within-LOW estimate,
together with an insignificant within-HIGH estimate, does not constitute
a significant LOW-minus-HIGH difference. The direct contrast, its uncertainty,
and any specified multiple-testing correction must be reported.

Economic evaluation raises further independent challenges.
More recently, \cite{JensenKellyMalamudPedersen2026} develop an
implementable efficient frontier in which machine-learning portfolio
weights are learned with trading costs explicitly incorporated into
the economic objective. Their research shows why simple forecast-then-
trade policies can fail to deliver net performance when strategies
rely on fleeting small-stock characteristics. This is a different
optimization objective from the current frozen decile portfolios,
but it establishes that cost-aware investment use of machine learning
is already a substantial published research area.


\cite{NovyMarxVelikov2016} document the economic consequences of trading
costs across anomalies, while \cite{DeMiguelEtAl2020} analyze the influence
of transaction costs when exploiting multiple stock characteristics.
\cite{HouXueZhang2020} emphasize the sensitivity of anomaly evidence
to replication design, microcaps and statistical testing.
\cite{HarveyLiuZhu2016} caution that multiple testing complicates
interpretation of large collections of predictive findings, and
\cite{McLeanPontiff2016} study attenuation of anomaly returns across
subsequent and post-publication periods.

The present portfolio analysis sorts next-month predicted returns into
deciles and evaluates $D10-D1$ with equal and lagged-market-equity value
weights. The cost-adjusted three-state results are based on portfolio
turnover from return-drifted holdings and a stipulated one-way cost of
50 basis points per traded dollar for both long and short legs.
This is a reproducible \emph{scenario}, not an estimate of actual
execution, borrowing, financing or market-impact costs. In particular,
positive modeled net returns are insufficient to establish feasible
capacity or implementable profitability.

\subsection{Potential research gap and conditional contribution}
\label{subsec:literature-gap}

Taken together, the preceding studies leave a possible, but not yet
demonstrated, narrower contribution: a unified assessment of how
cross-sectional \emph{ranking quality}, \emph{relative model accuracy},
and \emph{net portfolio value} behave across historically observable
VIX states when both models use an identical, transparent baseline
information set and direct regime comparisons are subjected to a common
inference protocol. The distinction between descriptive within-state
performance and formal between-state heterogeneity is central.

This framing is deliberately conditional on the results already obtained.
The production baseline exhibits little overall next-month stock-level
forecast-error improvement relative to its historical-mean benchmark.
The original binary classification shows positive rank information and
a positive modeled net equal-weighted spread for XGBoost in LOW months,
but direct LOW/HIGH differences remain imprecise.
In the three-state evaluation, XGBoost's equal-weighted cost-adjusted
mean spread is approximately $0.614\%$ per month for LOW and
$0.960\%$ for MIDDLE, compared with $-0.515\%$ for HIGH.
However, the three-state net pairwise differences do not survive the
specified Holm adjustment, and the corresponding value-weighted net
spreads are negative. Thus, the existing results do not establish a
statistically significant ordering of regimes or economically
implementable profits. Distinct claims about stronger benchmarks,
threshold robustness, microcap attribution and interpretation await
additional analysis.

A negative outcome in those subsequent exercises need not be concealed.
If state-conditional differences disappear under better benchmarks,
historically justified volatility thresholds, investability restrictions,
or robust sample splits, the result would delimit the reliability
of simple characteristic-based machine-learning strategies. Whether such
an outcome contributes beyond existing evidence remains a question
for the final competitor audit and faculty review.

\subsection{Research questions and transparent hypothesis framework}
\label{subsec:literature-hypotheses}

Let $S_t\in\{L,M,H\}$ denote the three-state historical VIX label of
formation month $t$. For forecasting model $m$ define
\[
IC_t^{m}=\mathrm{Spearman}_{i}
 \left(\widehat R_{i,t+1}^{m},R_{i,t+1}\right),
\]
and let $G_{w,t}^{m}$ and $N_{w,t}^{m}$ respectively denote gross and
50-basis-point modeled net returns on its $D10-D1$ strategy with
weighting convention $w\in\{EW,VW\}$. These monthly outcomes are
evaluated conditional on the formation-month state, not by
training models separately within states.

\paragraph{H1: Regime heterogeneity.}
For a declared outcome $Y_t^m$ (starting with monthly rank IC and
separately considering gross and net portfolio returns), the null is
\[
H_{0,1}:\quad
 E[Y_t^m\mid S_t=L]=E[Y_t^m\mid S_t=M]
 =E[Y_t^m\mid S_t=H].
\]
The two-sided alternative is that at least one conditional mean differs.
Joint time-series tests and multiplicity-controlled pairwise contrasts,
rather than isolated within-state significance, address this question.

\paragraph{H2: Relative nonlinear predictive advantage.}
Let
\[
\Delta_r^{IC}
 = E[IC_t^{XGB}-IC_t^{EN}\mid S_t=r].
\]
The state-heterogeneity null is
$H_{0,2}:\Delta_L^{IC}=\Delta_M^{IC}=\Delta_H^{IC}$,
with a corresponding null for aligned monthly forecast-loss
differences, where the loss sign is specified before estimation.
The relevant object is a \emph{paired} model difference, not
a comparison of two separately computed significance stars.

\paragraph{H3: Conditional economic value.}
The observed three-state net equal-weighted LOW XGBoost mean motivates
an examination of whether the conditional expected modeled net
$D10-D1$ return is positive and remains so under investability filters.
A positive sample mean, a statistically nonzero mean, factor-adjusted
performance and executable profit are not equivalent conclusions.
Tests already run on the current sample are retrospective evidence.

\paragraph{H4: Intermediate-volatility maximum (exploratory).}
A peak requires both
$E[Y_t\mid S_t=M]>E[Y_t\mid S_t=L]$ and
$E[Y_t\mid S_t=M]>E[Y_t\mid S_t=H]$
for a predeclared outcome. This possibility arose after the
three-state results were inspected. It is therefore exploratory, not
an ex ante confirmatory claim, and requires an appropriately joint
or simultaneous test plus additional independent evidence.

\paragraph{H5: Investability and stability.}
Future tests will evaluate whether any positive conditional net
spread remains after prespecified microcapitalization, liquidity,
borrow-cost and other implementability restrictions.
This is a future analysis plan, not established evidence from the
present cost scenario.

The full hypothesis registry, including estimands, dependence assumptions,
test families, decision rules, exploratory status and mappings to future
research milestones, is maintained in the project's Milestone 11
documentation. All hypotheses are documented after the M1--M10
results were observed. This section must not be interpreted as
retrospective preregistration.
\section{Data}

\subsection{CRSP monthly stock data}

The production data construction uses monthly U.S. common-stock data obtained through WRDS from the current CRSP CIZ monthly stock file, \texttt{crsp.msf\_v2}. The universe is restricted to U.S. ordinary common equity listed on the NYSE, AMEX, and NASDAQ using the available CIZ security classifications. The production research sample spans January 1990 through December 2025. Additional pre-sample CRSP history beginning in 1988 is retained during construction so that momentum and prior-December market equity are available at the start of the research sample.

Price is normalized to its absolute value. Market equity is constructed as price times shares outstanding. Because CRSP shares outstanding are reported in thousands, this monthly market-equity variable is measured in thousands of U.S. dollars. Log market equity is defined only for strictly positive market equity. Lagged market equity refers to the previous actual calendar month rather than simply the previous row.

A \$5 price screen is applied at feature month $t$ after the return history required for target construction has been assembled. This ordering ensures that the price realized at $t+1$ cannot determine whether the observation at $t$ survives the feature-month screen.

\subsection{Compustat fundamentals and book equity}

Annual accounting data are obtained from \texttt{comp.funda}. The production extraction uses standard industrial, consolidated, domestic-population annual data reported in U.S. dollars where the corresponding fields are available. Construction uses pre-sample Compustat history beginning in 1987.

Stockholders' equity is constructed using the hierarchy \texttt{SEQ}, then \texttt{CEQ + PSTK}, and then \texttt{AT - LT}. Preferred stock uses the hierarchy \texttt{PSTKRV}, then \texttt{PSTKL}, and then \texttt{PSTK}; when all preferred-stock fields are missing, preferred stock is explicitly set to zero. Deferred taxes and investment tax credit use \texttt{TXDITC}, with \texttt{TXDB + ITCB} as the fallback; when all such fields are missing, this component is explicitly set to zero.

Book equity is therefore
\begin{equation}
BE = SHE + DT - PS,
\end{equation}
where $SHE$ is stockholders' equity, $DT$ is deferred taxes and investment tax credit, and $PS$ is preferred stock. Nonpositive book equity is retained for data-lineage diagnostics but does not generate a valid baseline book-to-market characteristic.

\subsection{CRSP/Compustat linkage}

Compustat observations are linked to CRSP securities through the CRSP/Compustat Merged database. The implementation uses the link-history table \texttt{crsp.ccmxpf\_lnkhist}. The baseline retains link types \texttt{LC} and \texttt{LU} and primary indicators \texttt{P} and \texttt{C}. A link must be effective on the Compustat fiscal-year-end date. Open-ended link dates are handled as open intervals on the missing side.

Competing mappings are resolved deterministically before monthly assignment, prioritizing primary links and the accepted link types. The construction never merges Compustat to CRSP by ticker or company name.

\subsection{Accounting timing and book-to-market}

The baseline uses an explicit conservative June convention. For a Compustat fiscal-year end in calendar year $y-1$, the corresponding accounting characteristic first becomes eligible in June of year $y$ and remains active through May of year $y+1$. Accounting observations are therefore not forward-filled into months immediately following the fiscal-year-end date.

For characteristic year $y$, book equity is paired with CRSP market equity from December $y-1$. CRSP market equity is converted from thousands of dollars to millions of dollars before it is combined with Compustat book equity, which is reported in millions.

Where a firm has multiple eligible linked CRSP securities represented in the research universe, the December market-equity denominator is formed at the firm level by summing the required security-level market capitalizations. If a required linked security lacks its December market-equity observation, the denominator is treated as incomplete and book-to-market remains missing.

The baseline characteristic is
\begin{equation}
BM_{i,y} = \frac{BE_{i,y-1}}{ME^{Dec}_{i,y-1}}.
\end{equation}
A valid book-to-market observation requires strictly positive book equity and strictly positive, complete prior-December market equity. No winsorization or trimming is introduced during data construction.

\subsection{Momentum and target}

Momentum 12--2 compounds monthly stock returns from $t-12$ through $t-2$, excluding months $t$ and $t-1$. The target is the stock return in the next actual calendar month, $r_{i,t+1}$. Feature construction requires calendar adjacency; a gap in a security's history cannot cause a later observation to be treated as the next-month target or as part of a continuous momentum window.

The canonical model-ready panel retains an explicit formation date $t$ and realized-return date $t+1$ for every security-month. It also retains the accounting datadate, June eligibility date, fiscal year, prior-December market equity, and the contemporaneous VIX state so that timing can be audited without reconstructing predictors downstream. Raw predictor values are stored without full-sample standardization or implicit winsorization.

\subsection{VIX data and volatility regimes}

Market volatility is measured using the Cboe VIX series resolved in the production build from \texttt{cboe\_all.cboe}. Daily observations are converted to monthly frequency using the final observed VIX close in each calendar month.

For each month $t$, the volatility threshold is the expanding median of monthly VIX observations through $t$. A month is classified as HIGH when contemporaneous VIX exceeds this expanding median and LOW otherwise. Future volatility information therefore does not enter the regime assigned at the prediction date.

\subsection{Production panel coverage}

The validated production panel contains 2,140,091 stock-month observations across 18,898 PERMNOs from January 1990 through December 2025. Book-to-market is available for 1,818,006 stock-months, and 1,798,891 observations contain all three baseline predictors---log market equity, book-to-market, and momentum 12--2---before the sequential downstream eligibility filters. The final panel contains no duplicate PERMNO--date observations.

After sequentially requiring a valid next-month target, a feature-month price of at least \$5, valid size and momentum, a valid CCM match, positive book equity, and valid book-to-market, 1,340,823 stock-month observations remain in the complete baseline sample flow.

The empirical results reported elsewhere from the earlier preliminary analysis were generated before book-to-market was integrated and therefore use only size and momentum. They are retained as preliminary historical results rather than being reinterpreted as estimates from the completed three-predictor panel.

\section{Methods}

\subsection{Forecasting objective}

For stock $i$ in formation month $t$, the forecasting problem is
\begin{equation}
\widehat{r}_{i,t+1} = f\left(X_{i,t}\right),
\end{equation}
where $X_{i,t}$ contains characteristics observable at the formation date and $r_{i,t+1}$ is the realized next-calendar-month return.

The frozen baseline information set contains three predictors: log market equity, book-to-market, and momentum 12--2. The completed data-construction pipeline attaches these predictors to formation month $t$ using only information allowed by the CRSP market-data timing and the annual accounting convention described in the Data section.

\subsection{Accounting-information safeguards}

Annual book-to-market is assigned with an explicit characteristic-year key rather than unrestricted forward filling. A fiscal-year-end observation from calendar year $y-1$ first becomes eligible in June $y$ and remains active through May $y+1$. The market-equity denominator is prior-December firm market equity, converted to the same units as Compustat book equity.

CCM mappings must be effective on the accounting observation's fiscal-year-end date. Duplicate and competing mappings are resolved deterministically before monthly assignment. Validation rules reject accounting observations that would enter before their permitted availability date.

\subsection{Volatility regimes}

The original baseline volatility specification remains a binary formation-time VIX regime. Let $\widetilde{VIX}_t$ denote the expanding historical median computed through and including formation month $t$. The baseline state is HIGH when $VIX_t>\widetilde{VIX}_t$ and LOW otherwise. The state at $t$ classifies the portfolio whose realized stock returns occur at $t+1$.

As an additional conditional-evaluation specification, the analysis also constructs expanding historical VIX terciles using the same information-set convention. Let $Q_{0.33,t}$ and $Q_{0.67,t}$ denote the 33.33rd and 66.67th percentiles of VIX observations available through and including $t$. The implementation uses NumPy's deterministic linear quantile convention. The three states are
\begin{equation}
\begin{aligned}
\mathrm{LOW}_t: \quad &VIX_t\leq Q_{0.33,t}, \\
\mathrm{MIDDLE}_t: \quad &Q_{0.33,t}<VIX_t\leq Q_{0.67,t}, \\
\mathrm{HIGH}_t: \quad &VIX_t>Q_{0.67,t}.
\end{aligned}
\end{equation}
The thresholds evolve through time and never use VIX observations after the formation month. The three-state specification is an extension to, not a replacement for, the original expanding-median binary regime.

\subsection{Elastic Net}

Elastic Net provides the regularized linear benchmark. Predictor scaling is performed inside the model pipeline and is fit on historical training observations only. The selected production hyperparameters are $\alpha=0.01$ and $l_1$ ratio $=0.1$, with a maximum of 20,000 optimization iterations.

\subsection{XGBoost}

XGBoost provides the nonlinear forecasting model with squared-error loss and a fixed random seed. The selected production specification uses 200 trees, maximum depth 2, learning rate 0.01, column subsampling 0.7, row subsampling 1.0, histogram tree construction, and four fitting threads.

\subsection{Historical model selection}

Hyperparameters are selected once using only the 1990--1999 pre-OOS sample and five forward-chaining validation folds. To make the historical grid search computationally feasible while preserving calendar structure, tuning uses a deterministic cap of 500 observations per historical month. This sampling is confined to model selection; production OOS refits use the full eligible expanding historical sample. Selected hyperparameters are frozen before the January 2000 OOS period begins.

\subsection{Time-respecting out-of-sample design}

The empirical implementation uses an expanding historical training window and monthly refitting. For every prediction month, training feature dates are strictly earlier than the test month, and training outcomes must already be observable by the formation date. Concretely, a historical row formed at $s$ can enter the training sample at $t$ only when its recorded realized-return date satisfies $s+1\leq t$. No random train/test split is used.

The final production OOS panel contains 893,486 stock-month predictions across 311 formation months from January 2000 through November 2025. There are 138 HIGH-VIX and 173 LOW-VIX formation months.

\subsection{Forecast evaluation}

Forecast accuracy is summarized with mean squared error, root mean squared error, mean absolute error, prediction--realization correlation, and out-of-sample $R^2$.

The production OOS benchmark is the pooled expanding historical mean. At formation month $t$, every stock receives the mean of all historical next-month returns whose realization dates are observable on or before $t$:
\begin{equation}
R^2_{OOS}
=
1-
\frac{\sum_{i,t}(r_{i,t+1}-\widehat r_{i,t+1})^2}
{\sum_{i,t}(r_{i,t+1}-\bar r^{\,hist}_{t})^2}.
\end{equation}
This benchmark is fixed ex ante and differs from the stock-specific historical mean used in an earlier single-stock validation.

Cross-sectional ranking is evaluated each month using Spearman and Pearson correlations between predicted and realized returns. Mean monthly Spearman information coefficients are tested using Newey--West/HAC standard errors with six monthly lags. Direct LOW-minus-HIGH differences are estimated with a regression of monthly IC on an indicator for the LOW-VIX regime using the same HAC lag length.

\subsection{Economic-value evaluation}

Stocks are sorted each month into ten prediction portfolios using tie-preserving percentile ranks. D1 contains the lowest predicted returns and D10 the highest, and the long--short portfolio is $D10-D1$. Results are calculated with equal weights and with value weights based on lagged market equity. The lagged market-equity date is required to be the previous month-end.

Turnover is measured against the previous portfolio after drifting prior weights through realized returns. The frozen transaction-cost robustness specification charges 50 basis points one way per dollar of turnover. Gross and net monthly long--short returns are summarized overall and by VIX regime. Mean returns and LOW-minus-HIGH differences are evaluated using Newey--West/HAC inference with six monthly lags.


For the three-state volatility extension, stocks are ordered within each formation month and model by the frozen OOS forecast, with PERMNO used as a deterministic exact-tie breaker before forming ten approximately equal groups. The same cross-sectional sort is used in every volatility state; regimes classify monthly portfolio performance after sorting. This extension reports gross EW and lagged-ME VW results descriptively. Formal three-way HAC regime-difference tests and transaction-cost adjustments are reserved for later inference/robustness work rather than used to define or tune the three-state classification.

\section{Results}

\subsection{Production out-of-sample sample}

The production evaluation contains 893,486 stock-month predictions across 311 formation months from January 2000 through November 2025. The regime classification yields 138 HIGH-VIX months and 173 LOW-VIX months. Both Elastic Net and XGBoost produce complete finite predictions for every OOS observation.

\subsection{Forecast-error performance}

Table~\ref{tab:baseline-oos} reports return-level forecast performance relative to the pooled expanding historical-mean benchmark.

\begin{table}[htbp]
\centering
\caption{Production out-of-sample forecast performance}
\label{tab:baseline-oos}
\begin{tabular}{llrrrr}
\hline
Model & Regime & $N$ & RMSE & $R^2_{OOS}$ & Corr. \\
\hline
Elastic Net & Overall & 893486 & 0.13788 & -0.00027 & -0.02212 \\
Elastic Net & HIGH    & 398295 & 0.16132 & -0.00068 & -0.04289 \\
Elastic Net & LOW     & 495191 & 0.11562 &  0.00038 &  0.01351 \\
XGBoost     & Overall & 893486 & 0.13787 & -0.00023 & -0.00529 \\
XGBoost     & HIGH    & 398295 & 0.16135 & -0.00106 & -0.02999 \\
XGBoost     & LOW     & 495191 & 0.11558 &  0.00106 &  0.02914 \\
\hline
\end{tabular}
\end{table}

The return-level results are economically small. Neither model improves materially on the pooled historical-mean benchmark overall. Both models produce slightly positive OOS $R^2$ in LOW-VIX months and negative OOS $R^2$ in HIGH-VIX months, with XGBoost obtaining the largest LOW-VIX value at approximately 0.0011.

\subsection{Cross-sectional ranking}

Table~\ref{tab:rank-ic} reports monthly Spearman information coefficients with Newey--West/HAC inference.

\begin{table}[htbp]
\centering
\caption{Monthly cross-sectional Spearman information coefficients}
\label{tab:rank-ic}
\begin{tabular}{llrrrr}
\hline
Model & Regime & Months & Mean IC & HAC $t$ & $p$-value \\
\hline
Elastic Net & Overall & 311 &  0.00303 &  0.392 & 0.6952 \\
Elastic Net & HIGH    & 138 & -0.00516 & -0.403 & 0.6872 \\
Elastic Net & LOW     & 173 &  0.00957 &  1.160 & 0.2462 \\
XGBoost     & Overall & 311 &  0.00960 &  1.283 & 0.1996 \\
XGBoost     & HIGH    & 138 & -0.00178 & -0.143 & 0.8866 \\
XGBoost     & LOW     & 173 &  0.01867 &  2.479 & 0.0132 \\
\hline
\end{tabular}
\end{table}

The strongest ranking result is XGBoost in LOW-VIX months. Its mean monthly Spearman IC is 0.01867 and is statistically positive at the 5\% level. In HIGH-VIX months the XGBoost mean IC is close to zero and statistically insignificant. Elastic Net does not display statistically significant mean rank IC overall or within either regime.

The direct LOW-minus-HIGH rank-IC difference is 0.01473 for Elastic Net ($t=1.08$, $p=0.279$) and 0.02045 for XGBoost ($t=1.62$, $p=0.105$). The production evidence therefore shows statistically detectable XGBoost ranking ability within LOW-VIX states, but the direct regime contrast does not reach conventional significance.

\subsection{Three-state volatility-regime rank inference}

The three-state volatility extension divides formation months into LOW, MIDDLE, and HIGH states using expanding historical VIX tercile thresholds. The extension preserves the original binary regime, the frozen out-of-sample predictions, the predictor set, and the model-estimation procedure. Formal monthly inference uses Newey--West heteroskedasticity- and autocorrelation-consistent covariance with the repository's frozen six-month lag.

Table~\ref{tab:rank-ic-three-state-inference} reports within-regime monthly Spearman information coefficients and the corresponding HAC inference.

\begin{table}[htbp]
\centering
\caption{Three-state HAC inference for monthly Spearman information coefficients}
\label{tab:rank-ic-three-state-inference}
\small
\begin{tabular}{llrrrrrr}
\hline
Model & Regime & Mean IC & HAC SE & $t$ & $p$ & $N$ & Sig. \\
\hline
Elastic Net & LOW    &  0.00703 & 0.00868 &  0.81 & 0.4180 & 104 &  \\
Elastic Net & MIDDLE &  0.01257 & 0.01344 &  0.94 & 0.3494 & 108 &  \\
Elastic Net & HIGH   & -0.01157 & 0.01476 & -0.78 & 0.4330 &  99 &  \\
XGBoost     & LOW    &  0.01235 & 0.00760 &  1.63 & 0.1039 & 104 &  \\
XGBoost     & MIDDLE &  0.02129 & 0.01237 &  1.72 & 0.0852 & 108 & * \\
XGBoost     & HIGH   & -0.00605 & 0.01397 & -0.43 & 0.6648 &  99 &  \\
\hline
\end{tabular}
\begin{flushleft}
\footnotesize
Notes: IC is the monthly cross-sectional Spearman correlation between predicted and realized next-month returns. HAC inference uses the frozen six-month Newey--West lag. Significance stars denote $^{***}p<0.01$, $^{**}p<0.05$, and $^{*}p<0.10$.
\end{flushleft}
\end{table}

The three-state rank results display meaningful descriptive variation across volatility states, particularly for XGBoost. Its mean monthly Spearman IC is 0.01235 in LOW-volatility months, 0.02129 in MIDDLE-volatility months, and -0.00605 in HIGH-volatility months. The MIDDLE-regime mean is marginally positive at the 10\% level ($t=1.72$, $p=0.0852$). The LOW-regime mean is also positive but does not reach the 10\% threshold ($p=0.1039$), while the HIGH-regime mean is negative and statistically insignificant.

Within-regime significance does not by itself demonstrate that predictive performance differs statistically across regimes. To test regime differences directly, a LOW-reference regression is estimated for each model,

\begin{equation}
IC_t
=
\alpha
+
\beta_{\mathrm{MIDDLE}}
\mathbf{1}\{\mathrm{MIDDLE}_t\}
+
\beta_{\mathrm{HIGH}}
\mathbf{1}\{\mathrm{HIGH}_t\}
+
\varepsilon_t,
\end{equation}

with Newey--West/HAC covariance. The coefficient $\beta_{\mathrm{MIDDLE}}$ estimates MIDDLE minus LOW, $\beta_{\mathrm{HIGH}}$ estimates HIGH minus LOW, and HIGH minus MIDDLE is obtained as a covariance-consistent linear contrast. Holm adjustment is applied across the three pre-specified pairwise comparisons within each model family.

\begin{table}[htbp]
\centering
\caption{Three-state rank-IC regime-difference tests}
\label{tab:rank-ic-regime-differences}
\small
\begin{tabular}{llrrrrr}
\hline
Model & Comparison & Difference & HAC SE & $t$ & Raw $p$ & Holm $p$ \\
\hline
Elastic Net & MIDDLE - LOW  &  0.00554 & 0.01547 &  0.36 & 0.7201 & 0.7201 \\
Elastic Net & HIGH - LOW    & -0.01861 & 0.01576 & -1.18 & 0.2377 & 0.4754 \\
Elastic Net & HIGH - MIDDLE & -0.02415 & 0.01668 & -1.45 & 0.1476 & 0.4429 \\
XGBoost     & MIDDLE - LOW  &  0.00894 & 0.01434 &  0.62 & 0.5332 & 0.5332 \\
XGBoost     & HIGH - LOW    & -0.01841 & 0.01450 & -1.27 & 0.2043 & 0.4086 \\
XGBoost     & HIGH - MIDDLE & -0.02734 & 0.01537 & -1.78 & 0.0751 & 0.2254 \\
\hline
\end{tabular}
\begin{flushleft}
\footnotesize
Notes: Differences are covariance-consistent linear contrasts from LOW-reference three-state regressions. Newey--West/HAC covariance uses six monthly lags. Holm adjustment is applied across the three pre-specified comparisons within each model family.
\end{flushleft}
\end{table}

For XGBoost, the largest estimated rank-IC contrast is HIGH minus MIDDLE at -0.02734. The raw test gives $t=-1.78$ and $p=0.0751$, but the Holm-adjusted $p$-value is 0.2254. The HIGH-minus-LOW contrast is -0.01841 with a Holm-adjusted $p$-value of 0.4086. Elastic Net likewise shows no statistically significant pairwise contrast.

Accordingly, none of the pre-specified three-state rank-IC differences is statistically distinguishable from zero at the 5\% family-wise level after Holm adjustment. The three-state results therefore show descriptive and suggestive variation in cross-sectional ranking ability across volatility states, but do not establish statistically distinct rank predictability across LOW, MIDDLE, and HIGH regimes.

\subsection{Model interpretation}

The fitted Elastic Net coefficients indicate that book-to-market is the persistent linear predictor. Across monthly production refits, its standardized coefficient is nonzero in approximately 98.4\% of refits, whereas log market equity and momentum are nonzero much less frequently.

XGBoost uses all three predictors materially. Mean feature-importance shares across production refits are approximately 0.382 for book-to-market, 0.364 for momentum 12--2, and 0.254 for log market equity. These importance measures are descriptive and are not interpreted as causal or structural effects.

\subsection{Interpretation}

The production results distinguish return-level forecasting from cross-sectional ranking. OOS $R^2$ remains close to zero, yet XGBoost identifies statistically positive ranking ability in the original binary LOW-VIX specification. The three-state extension provides additional evidence that XGBoost rank performance is strongest descriptively in MIDDLE-volatility months and deteriorates in HIGH-volatility months. However, direct three-state contrasts do not remain statistically significant after Holm correction.

The combined evidence therefore suggests that the economically relevant predictive content is primarily ordinal across stocks and varies descriptively with volatility conditions, rather than representing a large reduction in individual-stock return squared error or a formally established difference among volatility regimes. The corresponding portfolio implications are evaluated in the next section.
\section{Economic Value}

The production economic-value analysis sorts stocks on out-of-sample predicted next-month returns within each formation month and model. D1 contains the lowest predictions and D10 the highest, and the long--short portfolio is $D10-D1$. The frozen three-state deciles use deterministic predicted-return order with PERMNO tie-breaking. Results are computed using equal weights and lagged-market-equity value weights.

The transaction-cost robustness specification charges 50 basis points one way per dollar of turnover. Turnover is measured relative to prior portfolio weights after they are drifted through realized returns.

\subsection{Gross and net long--short performance}

Table~\ref{tab:portfolio-performance} reports the principal production portfolio results. Net returns subtract the 50-basis-point turnover cost.

\begin{table}[htbp]
\centering
\caption{Production $D10-D1$ portfolio performance}
\label{tab:portfolio-performance}
{\small
\setlength{\tabcolsep}{3pt}
\begin{tabular}{lllrrrrr}
\hline
Model & Weight & Regime & Gross mean & Gross $t$ & Net mean & Net Sharpe & Net $t$ \\
\hline
Elastic Net & EW & Overall & 0.00342 & 1.084 &  0.00214 &  0.153 &  0.671 \\
Elastic Net & EW & HIGH    & 0.00096 & 0.169 & -0.00058 & -0.032 & -0.100 \\
Elastic Net & EW & LOW     & 0.00538 & 1.798 &  0.00430 &  0.441 &  1.451 \\
Elastic Net & VW & Overall & -0.00076 & -0.241 & -0.00198 & -0.122 & -0.623 \\
Elastic Net & VW & HIGH    & -0.00191 & -0.341 & -0.00358 & -0.170 & -0.630 \\
Elastic Net & VW & LOW     & 0.00016 & 0.049 & -0.00071 & -0.064 & -0.221 \\
XGBoost     & EW & Overall & 0.00617 & 1.920 &  0.00373 &  0.267 &  1.157 \\
XGBoost     & EW & HIGH    & 0.00176 & 0.304 & -0.00098 & -0.054 & -0.168 \\
XGBoost     & EW & LOW     & 0.00969 & 3.605 &  0.00748 &  0.796 &  2.815 \\
XGBoost     & VW & Overall & 0.00037 & 0.097 & -0.00248 & -0.144 & -0.639 \\
XGBoost     & VW & HIGH    & -0.00004 & -0.007 & -0.00311 & -0.145 & -0.454 \\
XGBoost     & VW & LOW     & 0.00071 & 0.219 & -0.00197 & -0.153 & -0.610 \\
\hline
\end{tabular}
}
\end{table}

The strongest economic-value result is the XGBoost equal-weighted LOW-VIX portfolio. Its gross mean return is 0.969\% per month with HAC $t=3.61$. After the frozen transaction-cost adjustment, the mean remains 0.748\% per month, corresponding to an annualized arithmetic mean of approximately 8.98\%, a net Sharpe ratio of 0.80, HAC $t=2.82$, and $p=0.0049$.

The corresponding HIGH-VIX equal-weighted XGBoost portfolio has a net mean of approximately -0.098\% per month and no evidence of statistical significance. Elastic Net's LOW-VIX equal-weighted spread is positive but weaker: 0.430\% per month net, with HAC $t=1.45$ and $p=0.147$.

\subsection{Value weighting}

The XGBoost signal does not survive value weighting. In LOW-VIX months the gross value-weighted spread is only 0.071\% per month, and the net spread after turnover costs is -0.197\% per month. Elastic Net value-weighted portfolios are similarly weak.

The contrast between equal- and value-weighted results indicates that the production economic signal is concentrated away from the largest stocks. This is an economically important qualification because smaller-stock strategies can face higher implementation frictions than the baseline transaction-cost assumption captures.

\subsection{Direct binary regime tests}

Table~\ref{tab:portfolio-regime-tests} reports LOW-minus-HIGH HAC tests for the original binary volatility specification.

\begin{table}[htbp]
\centering
\caption{LOW-minus-HIGH long--short portfolio differences}
\label{tab:portfolio-regime-tests}
\begin{tabular}{lllrrr}
\hline
Model & Weight & Return & LOW--HIGH & HAC $t$ & $p$-value \\
\hline
Elastic Net & EW & Gross & 0.00443 & 0.807 & 0.4196 \\
Elastic Net & EW & Net   & 0.00488 & 0.874 & 0.3819 \\
Elastic Net & VW & Gross & 0.00206 & 0.335 & 0.7377 \\
Elastic Net & VW & Net   & 0.00287 & 0.460 & 0.6456 \\
XGBoost     & EW & Gross & 0.00793 & 1.432 & 0.1521 \\
XGBoost     & EW & Net   & 0.00846 & 1.521 & 0.1282 \\
XGBoost     & VW & Gross & 0.00075 & 0.110 & 0.9121 \\
XGBoost     & VW & Net   & 0.00115 & 0.168 & 0.8669 \\
\hline
\end{tabular}
\end{table}

Although the XGBoost equal-weighted LOW-VIX strategy is statistically positive within LOW months, the direct LOW-minus-HIGH portfolio difference is not conventionally significant. The correct interpretation is therefore that economic predictability is concentrated in LOW-VIX states in this baseline specification, not that the LOW regime is formally proven to outperform the HIGH regime.

\subsection{Three-state volatility-regime portfolio inference}

The three-state extension separates formation months into LOW, MIDDLE, and HIGH volatility states using expanding historical VIX tercile thresholds. This extension evaluates the already-frozen out-of-sample predictions and does not retrain the models, change forecasts, modify portfolio assignments, alter portfolio weights, or redefine the original binary regime.

Formal inference is conducted on the monthly gross $D10-D1$ return series. Newey--West/HAC covariance uses the same frozen six-month lag as the existing production inference.

Table~\ref{tab:three-state-long-short-inference} reports within-regime results.

\begin{table}[htbp]
\centering
\caption{HAC inference for monthly $D10-D1$ returns across three-state VIX regimes}
\label{tab:three-state-long-short-inference}
\small
\begin{tabular}{lllrrrrrr}
\hline
Model & Weight & Regime & Mean & HAC SE & $t$ & $p$ & $N$ & Sig. \\
\hline
Elastic Net & EW & LOW    &  0.00486 & 0.00281 &  1.73 & 0.0836 & 104 & * \\
Elastic Net & EW & MIDDLE &  0.00670 & 0.00545 &  1.23 & 0.2190 & 108 &  \\
Elastic Net & EW & HIGH   & -0.00174 & 0.00664 & -0.26 & 0.7935 &  99 &  \\
Elastic Net & VW & LOW    &  0.00344 & 0.00308 &  1.12 & 0.2637 & 104 &  \\
Elastic Net & VW & MIDDLE & -0.00314 & 0.00466 & -0.67 & 0.5011 & 108 &  \\
Elastic Net & VW & HIGH   & -0.00250 & 0.00719 & -0.35 & 0.7276 &  99 &  \\
XGBoost     & EW & LOW    &  0.00813 & 0.00253 &  3.21 & 0.0013 & 104 & *** \\
XGBoost     & EW & MIDDLE &  0.01201 & 0.00536 &  2.24 & 0.0249 & 108 & ** \\
XGBoost     & EW & HIGH   & -0.00232 & 0.00661 & -0.35 & 0.7262 &  99 &  \\
XGBoost     & VW & LOW    &  0.00171 & 0.00317 &  0.54 & 0.5899 & 104 &  \\
XGBoost     & VW & MIDDLE &  0.00089 & 0.00604 &  0.15 & 0.8831 & 108 &  \\
XGBoost     & VW & HIGH   & -0.00149 & 0.00840 & -0.18 & 0.8594 &  99 &  \\
\hline
\end{tabular}
\begin{flushleft}
\footnotesize
Notes: Returns are decimal monthly returns. HAC inference uses Newey--West covariance with the repository's frozen six-month lag. Significance stars denote $^{***}p<0.01$, $^{**}p<0.05$, and $^{*}p<0.10$.
\end{flushleft}
\end{table}

The clearest three-state result appears in the equal-weighted XGBoost portfolios. Mean monthly $D10-D1$ returns are 0.813\% in LOW-volatility months ($t=3.21$, $p=0.0013$), 1.201\% in MIDDLE-volatility months ($t=2.24$, $p=0.0249$), and -0.232\% in HIGH-volatility months ($t=-0.35$, $p=0.7262$). Thus, the spread is statistically positive within both LOW and MIDDLE volatility states but not within HIGH states.

Elastic Net equal-weighted performance is weaker. Its LOW-regime mean is 0.486\% per month and is marginally significant at the 10\% level ($t=1.73$, $p=0.0836$), while MIDDLE and HIGH results are statistically insignificant. Neither model exhibits statistically significant three-state spreads under value weighting.

Within-regime significance, however, does not establish that the mean spread differs statistically across volatility states. Direct regime differences are therefore estimated using a LOW-reference regression,

\begin{equation}
R^{LS}_t
=
\alpha
+
\beta_{\mathrm{MIDDLE}}
\mathbf{1}\{\mathrm{MIDDLE}_t\}
+
\beta_{\mathrm{HIGH}}
\mathbf{1}\{\mathrm{HIGH}_t\}
+
\varepsilon_t,
\end{equation}

where $\beta_{\mathrm{MIDDLE}}$ identifies MIDDLE minus LOW, $\beta_{\mathrm{HIGH}}$ identifies HIGH minus LOW, and HIGH minus MIDDLE is obtained using the full HAC covariance matrix. Holm adjustment is applied across the three pre-specified pairwise comparisons separately within each model-by-weighting family.

\begin{table}[htbp]
\centering
\caption{Pairwise three-state volatility-regime differences in $D10-D1$ returns}
\label{tab:three-state-long-short-differences}
\small
\begin{tabular}{lllrrrrr}
\hline
Model & Weight & Comparison & Difference & HAC SE & $t$ & Raw $p$ & Holm $p$ \\
\hline
Elastic Net & EW & MIDDLE - LOW  &  0.00184 & 0.00546 &  0.34 & 0.7364 & 0.7689 \\
Elastic Net & EW & HIGH - LOW    & -0.00660 & 0.00685 & -0.96 & 0.3352 & 0.7689 \\
Elastic Net & EW & HIGH - MIDDLE & -0.00844 & 0.00744 & -1.14 & 0.2563 & 0.7689 \\
Elastic Net & VW & MIDDLE - LOW  & -0.00658 & 0.00525 & -1.25 & 0.2104 & 0.6311 \\
Elastic Net & VW & HIGH - LOW    & -0.00595 & 0.00784 & -0.76 & 0.4480 & 0.8960 \\
Elastic Net & VW & HIGH - MIDDLE &  0.00063 & 0.00875 &  0.07 & 0.9424 & 0.9424 \\
XGBoost     & EW & MIDDLE - LOW  &  0.00388 & 0.00547 &  0.71 & 0.4789 & 0.4789 \\
XGBoost     & EW & HIGH - LOW    & -0.01045 & 0.00689 & -1.52 & 0.1292 & 0.2584 \\
XGBoost     & EW & HIGH - MIDDLE & -0.01432 & 0.00736 & -1.95 & 0.0516 & 0.1548 \\
XGBoost     & VW & MIDDLE - LOW  & -0.00082 & 0.00660 & -0.12 & 0.9012 & 1.0000 \\
XGBoost     & VW & HIGH - LOW    & -0.00320 & 0.00861 & -0.37 & 0.7107 & 1.0000 \\
XGBoost     & VW & HIGH - MIDDLE & -0.00238 & 0.00962 & -0.25 & 0.8050 & 1.0000 \\
\hline
\end{tabular}
\begin{flushleft}
\footnotesize
Notes: Differences are covariance-consistent linear contrasts from LOW-reference three-state regressions. Newey--West/HAC covariance uses six monthly lags. Holm adjustment is applied within each pre-specified three-comparison model-by-weighting family. Benjamini--Hochberg FDR-adjusted values are retained in the machine-readable research outputs.
\end{flushleft}
\end{table}

For equal-weighted XGBoost, the estimated HIGH-minus-LOW difference is -1.045 percentage points per month ($t=-1.52$, raw $p=0.1292$). The largest contrast is HIGH minus MIDDLE at -1.432 percentage points per month, with $t=-1.95$ and raw $p=0.0516$. After Holm correction, however, the corresponding $p$-values are 0.2584 and 0.1548.

Consequently, none of the three pairwise XGBoost equal-weighted regime contrasts is statistically significant at the 5\% family-wise level. No Elastic Net or value-weighted pairwise comparison survives the Holm adjustment either.

The correct interpretation is therefore narrower than the raw within-regime estimates alone might suggest. XGBoost equal-weighted portfolios display economically large and statistically positive spreads in LOW and MIDDLE volatility states and weak performance in HIGH-volatility states, but the formal direct tests do not establish statistically distinct three-state regime effects after correction for multiple comparisons.

\subsection{Economic interpretation}

The portfolio results reinforce the cross-sectional ranking evidence while adding an important weighting dimension. In the original binary specification, the nonlinear model produces economically meaningful equal-weighted performance in LOW-VIX months, while the direct LOW-versus-HIGH contrast remains imprecisely estimated.

The three-state extension adds a more detailed pattern. Equal-weighted XGBoost spreads are positive and statistically significant within LOW and MIDDLE states and negative but statistically insignificant within HIGH states. The largest estimated deterioration occurs between MIDDLE and HIGH volatility. However, the formal pairwise regime tests do not remain significant after Holm multiple-testing adjustment.

The disappearance of the signal under value weighting remains an important qualification. It is consistent with predictive performance being stronger away from the largest firms, although the present analysis does not by itself identify firm size as the causal source of the weighting difference. Smaller-stock strategies may also face implementation frictions beyond those captured by the baseline transaction-cost specification.

Overall, the economic-value evidence supports meaningful state-dependent descriptive variation in the frozen out-of-sample prediction signal, especially for XGBoost under equal weighting. It does not establish that LOW, MIDDLE, and HIGH volatility regimes have statistically distinct expected long--short returns after family-wise multiple-testing correction. This distinction motivates subsequent robustness analysis rather than stronger regime-effect claims.

% Validated transaction-cost analysis
\subsection{Trading intensity and cost-adjusted three-state results}

To isolate implementation frictions without re-estimating forecasts, the turnover analysis
consumes the frozen month-by-month M8 assignments. The equal-weighted target is $1/N_{d,t}$;
value-weighted targets use the previously validated positive lagged market equity of stocks
in each decile. The current pre-trade portfolio weights are last month's targets drifted
through the return realized by the current formation date:
\begin{equation}
 \widetilde{w}_{i,t}=\frac{w_{i,t-1}(1+R_{i,t})}{\sum_jw_{j,t-1}(1+R_{j,t})},
 \qquad
 T_{d,t}=\frac12\sum_i |w_{i,t}-\widetilde{w}_{i,t}|.
\end{equation}
The first formation from cash has $T_{d,t}=1$. The sum runs over the union of
old and new holdings, including full liquidations and new purchases. Returns
$R_{i,t}$ in the drift are the prior formation's observed holding-period returns;
formation-date forecasts do not depend on their own not-yet-realized $t+1$ returns.

The frozen convention is one dollar long in D10 and one dollar short in D1
(two dollars gross and zero dollars net exposure). With one-way cost $c$ the net spread is
\begin{equation}
 R^{LS,\mathrm{net}}_{t+1}=R^{D10}_{t+1}-R^{D1}_{t+1}
 -c\left(T_{D10,t}+T_{D1,t}\right),\qquad c\in\{0,0.005\}.
\end{equation}
The 0-basis-point analysis must exactly reproduce the M8 gross series.
Both long and short legs incur rebalancing fees; the assumed 50-basis-point
rate is stylized and omits financing, borrow fees and nonlinear market impact.
All net-return HAC tests use the original six-lag M9 Newey--West covariance
and the same LOW-reference pairwise contrasts and Holm testing families.

\begin{table}[htbp]
\centering
\caption{Monthly gross versus 50-bp net $D10-D1$ by VIX regime}
\label{tab:transaction-costs}
\small
\begin{tabular}{lllrrrrrrr}
\hline
Model & Weight & Regime & Turnover & Gross & Net & HAC SE & $t$ & $p$ & $N$ \\
\hline
Elastic Net & EW & OVERALL & 0.256 & 0.0034 & 0.0021 & 0.0032 & 0.66 & 0.5067 & 311 \\
Elastic Net & EW & HIGH & 0.334 & -0.0017 & -0.0034 & 0.0067 & -0.51 & 0.6130 & 99 \\
Elastic Net & EW & LOW & 0.176 & 0.0049 & 0.0040 & 0.0028 & 1.43 & 0.1533 & 104 \\
Elastic Net & EW & MIDDLE & 0.263 & 0.0067 & 0.0054 & 0.0054 & 0.99 & 0.3201 & 108 \\
Elastic Net & VW & OVERALL & 0.243 & -0.0007 & -0.0020 & 0.0032 & -0.61 & 0.5403 & 311 \\
Elastic Net & VW & HIGH & 0.371 & -0.0025 & -0.0044 & 0.0073 & -0.60 & 0.5506 & 99 \\
Elastic Net & VW & LOW & 0.132 & 0.0034 & 0.0028 & 0.0031 & 0.91 & 0.3615 & 104 \\
Elastic Net & VW & MIDDLE & 0.234 & -0.0031 & -0.0043 & 0.0046 & -0.93 & 0.3503 & 108 \\
XGBoost & EW & OVERALL & 0.481 & 0.0062 & 0.0037 & 0.0032 & 1.16 & 0.2464 & 311 \\
XGBoost & EW & HIGH & 0.567 & -0.0023 & -0.0052 & 0.0067 & -0.77 & 0.4388 & 99 \\
XGBoost & EW & LOW & 0.400 & 0.0081 & 0.0061 & 0.0025 & 2.45 & 0.0143 & 104 \\
XGBoost & EW & MIDDLE & 0.482 & 0.0120 & 0.0096 & 0.0053 & 1.80 & 0.0714 & 108 \\
XGBoost & VW & OVERALL & 0.572 & 0.0004 & -0.0025 & 0.0039 & -0.64 & 0.5253 & 311 \\
XGBoost & VW & HIGH & 0.646 & -0.0015 & -0.0047 & 0.0085 & -0.56 & 0.5783 & 99 \\
XGBoost & VW & LOW & 0.505 & 0.0017 & -0.0008 & 0.0032 & -0.26 & 0.7964 & 104 \\
XGBoost & VW & MIDDLE & 0.568 & 0.0009 & -0.0020 & 0.0061 & -0.32 & 0.7476 & 108 \\
\hline
\end{tabular}
\begin{flushleft}
\footnotesize Notes: 50 basis points per one-sided turnover dollar; turnover is $D10+D1$ rebalancing turnover (initial leg turnover=1). Drift uses prior holding-period returns realized by each formation date. HAC uses six monthly lags. LOW/MIDDLE/HIGH use expanding formation-time VIX terciles.
\end{flushleft}
\end{table}


\section{Conclusion}

The research currently establishes a working, reproducible empirical foundation for studying regime-dependent stock-return predictability. The pipeline pulls real CRSP monthly data through WRDS, constructs predictors and next-month targets using explicit calendar timing, labels volatility regimes from an expanding historical VIX median, estimates Elastic Net and XGBoost under forward-looking restrictions, stores genuine out-of-sample predictions, and evaluates cross-sectional portfolio sorts with explicit safeguards against degenerate prediction cross-sections.

The preliminary real-data analysis covers 83 out-of-sample formation months from January 2019 through November 2025 using the partial size-and-momentum predictor set. The earlier AAPL validation shows that neither model outperforms the stock-specific expanding historical-mean benchmark, although XGBoost produces modestly lower forecast errors than Elastic Net. The subsequent cross-sectional analysis sharpens the interpretation. The tuned Elastic Net assigns one identical forecast to all eligible stocks in every OOS month and therefore contains no usable cross-sectional ranking signal in the current specification. Those month-model groups are excluded from decile formation rather than converted into arbitrary portfolios through identifier-based tie-breaking.

XGBoost produces non-degenerate prediction cross-sections in all 83 OOS months, but its mean monthly Spearman rank correlation with realized next-month returns is slightly negative. Equal-weighted D10-minus-D1 returns are economically close to zero overall. The value-weighted spread is positive overall and is concentrated in LOW-VIX months, while the HIGH-VIX value-weighted spread is essentially zero. Thus the current evidence does not support a simple claim that return predictability is stronger during high-volatility periods. Instead, it indicates that any preliminary economic-value signal is sensitive to portfolio weighting and appears concentrated among larger stocks during low-volatility months.

These findings remain descriptive rather than final empirical conclusions. The next stages add book-to-market through Compustat/CCM, expand the predictor set and production sample, finalize delisting-return treatment, implement turnover and the frozen 0/50-basis-point transaction-cost robustness specification, and add formal HAC/Newey--West inference and additional robustness analysis. The central research question therefore remains open, but the project has progressed from software and single-stock validation to a validated cross-sectional economic-value experiment with clearly identified model limitations.


\appendix
\section{Research Terminology and Methodology Reference}
\label{app:terminology}

This appendix records the principal terminology and conventions used in the production analysis. It is an academically focused subset of the repository's living technical reference, \texttt{docs/research\_glossary.md}. When an implementation detail is material to interpretation, the definitions below follow the validated production code and configuration rather than a generic textbook convention.

\subsection{Data sources and identifiers}

\begin{description}
\item[WRDS.] Wharton Research Data Services, the access platform used for licensed CRSP, Compustat, CCM, and the preferred Cboe VIX source.
\item[CRSP.] Center for Research in Security Prices. The production pipeline uses monthly U.S. security data from the current CIZ interface, with \texttt{crsp.msf\_v2} as the validated production source.
\item[CIZ.] The current integrated CRSP stock-data format used by the production source resolver. The production CIZ path identifies eligible ordinary common equity through the available security classifications rather than relying on legacy share codes alone.
\item[Compustat.] Source of annual firm accounting fundamentals; the validated production source is \texttt{comp.funda}.
\item[CCM.] CRSP/Compustat Merged link history. The production implementation uses \texttt{crsp.ccmxpf\_lnkhist} to map Compustat firms to CRSP securities subject to effective link dates.
\item[PERMNO.] CRSP permanent security identifier. It is the authoritative longitudinal security key in the monthly panel; ticker is retained for readability rather than used as the primary join key.
\item[GVKEY.] Compustat firm identifier used for annual fundamentals and CCM linkage.
\end{description}

\subsection{Returns, formation, and information timing}

Let $i$ index securities and let $t$ denote a month-end formation date. The forecasting target is the next actual calendar-month return $R_{i,t+1}$. The canonical panel stores the formation date and the realized-return date separately, and the latter must equal the next calendar month-end. A later available observation after a calendar gap is not substituted for $t+1$.

A historical observation may enter a model fitted for formation month $t$ only when its feature date is strictly earlier than $t$ and its target has already become observable by $t$. Thus, if a historical row was formed at month $s$, the row is usable only when its recorded realized-return date is no later than the current prediction date. This is the central target-observability safeguard against look-ahead bias.

\subsection{Stock characteristics}

\paragraph{Market equity (ME).}
For a CRSP security-month,
\begin{equation}
ME_{i,t}=|P_{i,t}|\times SHROUT_{i,t}.
\end{equation}
Because CRSP shares outstanding are measured in thousands, the core monthly market-equity variable is measured in thousands of U.S. dollars. Log size is
\begin{equation}
\texttt{log\_me}_{i,t}=\log(ME_{i,t}),
\end{equation}
and is defined only for strictly positive market equity.

The variable \texttt{me\_lag} is the previous actual calendar month's market equity. It is constructed on a calendar-complete monthly series, so a missing calendar month produces a missing lag rather than converting the previous observed row into a false one-month lag. Value-weighted portfolios use positive finite \texttt{me\_lag} measured at $t-1$.

\paragraph{Book equity (BE).}
Stockholders' equity follows the hierarchy \texttt{SEQ}, then \texttt{CEQ + PSTK}, then \texttt{AT - LT}. Preferred stock uses \texttt{PSTKRV}, then \texttt{PSTKL}, then \texttt{PSTK}, with zero used only when all preferred-stock fields are missing. Deferred taxes and investment tax credit use \texttt{TXDITC}, then \texttt{TXDB + ITCB}, with zero when the component is otherwise unavailable. The baseline book-equity definition is
\begin{equation}
BE = SHE + DT - PS,
\end{equation}
where $SHE$ is stockholders' equity, $DT$ is deferred taxes and investment tax credit, and $PS$ is preferred stock. Nonpositive BE does not generate a valid baseline book-to-market observation.

\paragraph{Book-to-market (B/M).}
For characteristic year $y$,
\begin{equation}
BM_{i,y}=\frac{BE_{i,y-1}}{ME^{Dec}_{i,y-1}}.
\end{equation}
The denominator is complete prior-December firm market equity, converted from CRSP dollar-thousands to millions to match Compustat units. When a firm has multiple valid linked CRSP securities represented in the research universe, their required December market equity is summed at the firm level. If a required linked security lacks December ME, B/M is left missing rather than constructed from an incomplete denominator.

The accounting characteristic associated with a fiscal-year end in calendar year $y-1$ first becomes eligible in June $y$ and remains assigned through May $y+1$. This conservative June convention is enforced explicitly rather than through unrestricted forward filling.

\paragraph{Momentum 12--2.}
The momentum characteristic compounds monthly returns from $t-12$ through $t-2$:
\begin{equation}
MOM_{i,t}^{12-2}=\prod_{j=2}^{12}(1+R_{i,t-j})-1.
\end{equation}
Month $t-1$ and the formation month $t$ are excluded. Eleven calendar-contiguous historical return observations are required; missing calendar months break the window.

\subsection{VIX volatility regimes}

Daily Cboe VIX observations are converted to monthly frequency using the final observed VIX close in each calendar month. The original binary baseline uses the expanding median through and including formation month $t$. It is HIGH when $VIX_t>\widetilde{VIX}_t$ and LOW otherwise; the historical \texttt{regime} label remains this baseline definition.

The three-state extension adds expanding historical tercile thresholds without altering the binary labels. Let $Q_{0.33,t}$ and $Q_{0.67,t}$ denote the 33.33rd and 66.67th percentiles of VIX observations available through and including $t$, computed with the deterministic linear quantile convention. Then
\begin{equation}
\begin{aligned}
\mathrm{LOW}_t: \quad &VIX_t\leq Q_{0.33,t}, \\
\mathrm{MIDDLE}_t: \quad &Q_{0.33,t}<VIX_t\leq Q_{0.67,t}, \\
\mathrm{HIGH}_t: \quad &VIX_t>Q_{0.67,t}.
\end{aligned}
\end{equation}
These thresholds evolve through time; future VIX observations and portfolio outcomes never enter their construction. The state at formation month $t$ classifies performance realized at $t+1$. The variables \texttt{vix\_expanding\_q33}, \texttt{vix\_expanding\_q67}, \texttt{regime\_binary}, and \texttt{regime\_3state} make the two definitions explicit.

\subsection{Forecasting models and model selection}

\paragraph{Elastic Net.}
Elastic Net is the regularized linear benchmark combining $L_1$ and $L_2$ coefficient penalties. Its preprocessing pipeline contains \texttt{StandardScaler}, and the scaler is fitted on the historical training sample only. The selected production specification is $\alpha=0.01$, $l_1$ ratio $=0.1$, and a maximum of 20,000 optimization iterations. The model stage does not impute missing predictors; production estimation is complete-case for the three frozen predictors.

\paragraph{XGBoost.}
Extreme Gradient Boosting is the nonlinear boosted-tree model. The selected production specification uses 200 trees, maximum depth 2, learning rate 0.01, row subsampling 1.0, column subsampling 0.7, histogram tree construction, and four fitting threads, with squared-error loss and a fixed random seed.

\paragraph{Hyperparameters and forward-chaining validation.}
Model parameters are quantities estimated during fitting; hyperparameters control the fitting procedure or model architecture and are selected before the production OOS period. Production hyperparameters are selected once from the 1990--1999 pre-OOS history with five chronological forward-chaining folds. Validation training dates precede the validation month, and training targets must be observable by the corresponding validation date. A deterministic cap of 500 observations per historical month is used only for the grid search; full eligible historical samples are used for production refits. Selected hyperparameters are then frozen.

\subsection{Out-of-sample evaluation}

The production design uses an expanding historical training window and monthly refitting. For every OOS formation month, both models are estimated on all eligible historical observations permitted by the information-timing rules and then predict the current eligible cross-section.

The headline benchmark is the pooled expanding historical mean. At formation month $t$, every stock receives the mean of all next-month returns whose realization dates are observable by $t$. The production OOS coefficient of determination is
\begin{equation}
R^2_{OOS}
=1-
\frac{\sum_{i,t}(R_{i,t+1}-\widehat R_{i,t+1})^2}
{\sum_{i,t}(R_{i,t+1}-\bar R^{\,hist}_{t})^2}.
\end{equation}
A negative $R^2_{OOS}$ means that the model has larger squared forecast error than this historical benchmark over the evaluated observations; it does not mean ``negative accuracy.'' Mean squared error (MSE), root mean squared error (RMSE), mean absolute error (MAE), and prediction--realization correlation are also reported.

\subsection{Cross-sectional ranking and portfolios}

The monthly rank information coefficient (IC) is the cross-sectional Spearman correlation between predicted next-month returns formed at $t$ and realized returns at $t+1$. It is computed within each formation month and then summarized through time. Pearson cross-sectional correlation is retained as an additional diagnostic.

For economic-value evaluation, stocks are assigned to ten prediction portfolios within each month and model using tie-preserving percentile ranks. D1 contains the lowest predicted returns and D10 the highest. The long--short return is
\begin{equation}
R^{LS}_{t+1}=R^{D10}_{t+1}-R^{D1}_{t+1}.
\end{equation}
Equal-weighted (EW) portfolios assign equal target weights within each leg. Value-weighted (VW) portfolios use positive finite previous-calendar-month market equity, \texttt{me\_lag}, as the weighting variable.

Turnover is measured after drifting the prior month's portfolio weights through realized security returns. One-way turnover on a leg is one half of the $L_1$ distance between the new target weights and the return-drifted previous weights; a newly initiated leg has turnover one. Total long--short turnover is the sum across the long and short legs.

The frozen transaction-cost robustness specification charges 50 basis points one way per dollar of turnover. Therefore,
\begin{equation}
R^{net}_{t+1}=R^{gross}_{t+1}-0.005\times Turnover_{t}.
\end{equation}

\subsection{Performance and statistical inference}

Mean monthly long--short returns are reported gross and net of the frozen transaction-cost specification. The arithmetic annualized mean is $12\bar R$, annualized volatility is $\sqrt{12}\,s_R$, and the reported descriptive annualized Sharpe ratio is $\sqrt{12}\,\bar R/s_R$ under the implementation's zero-risk-free-rate convention.

Formal inference for monthly portfolio-return and rank-IC series uses heteroskedasticity- and autocorrelation-consistent (HAC) standard errors with Bartlett/Newey--West weights and six monthly lags in the frozen baseline. LOW-minus-HIGH differences are estimated directly rather than inferred by comparing two separately significant within-regime estimates.

\subsection{Data quality and leakage safeguards}

The canonical production panel has a unique security-month key defined by PERMNO and formation date. Complete-case eligibility requires the three baseline predictors, the next-month target, the VIX regime, and required timing fields after the feature-month price screen. The production pipeline does not apply baseline winsorization or full-sample predictor standardization. Elastic Net scaling is training-only; XGBoost uses the raw complete-case predictor values supplied to its estimator.

The principal safeguards against future information are: calendar-safe $t/t+1$ target construction; conservative June accounting assignment; effective-dated CCM links; prior-December B/M market equity; an expanding VIX threshold using information only through formation; pre-OOS forward-chaining hyperparameter selection; training-only preprocessing; and explicit target-observability requirements in every model refit.

\subsection{Historical and deferred terminology}

The earlier preliminary specification used size and momentum only. Its single-stock and reduced-cross-section outputs remain historical validation artifacts and are not reinterpreted as estimates from the current three-predictor baseline. The current production baseline includes book-to-market and the completed economic-value stage, including turnover, the 50-basis-point transaction-cost robustness specification, and HAC inference.

Fama--MacBeth regressions and SHAP explanations are not part of the current production baseline described in this manuscript appendix. If later research stages introduce them, their definitions and implementation status must be added to the living technical reference before that stage is considered complete.


\clearpage
\begingroup
\singlespacing
\setlength{\bibsep}{\baselineskip}
\bibliographystyle{chicago}
\bibliography{references}
\endgroup
\end{document}
